\makeabstractSHORT
{ % #1: Title
Generalized Bracket Rules for $B(\infty)$ Crystals
}
{ % #2: Authorship
Phillip Antis\textsuperscript{1},
Kiran Bishop\textsuperscript{2},
Timothy Carlin-Burns\textsuperscript{3},
Maren Forrest\textsuperscript{4},
Michael Myers\textsuperscript{5},
Verity Nwabuisi\textsuperscript{6}
}
{ % #3: Affiliations (use \\ to start a newline)
\textsuperscript{1} Department of Mathematics, Millsaps College, Jackson, MS\\
\textsuperscript{2} Department of Mathematics, Harvard University, Cambridge, MA\\
\textsuperscript{3} Department of Mathematical Sciences, Carnegie Mellon University, Pittsburgh, PA\\
\textsuperscript{4} Department of Mathematics, Amherst College, Amherst, MA\\
\textsuperscript{5} MSCS Department, Macalester College, St. Paul, MN\\
\textsuperscript{6} School of Mathematics and Natural Sciences, The University of Southern Mississippi, Hattiesburg, MS
}
{ % #4: Abstract
\textbf{Abstract:}
The crystal $B(\infty)$ is an infinite-dimensional crystal basis which is associated with the lower half of a quantum group. It acts as a combinatorial model, encoding information about its semisimple Lie Algebra and finite-dimensional representations. 

Lusztig, Berenstein, and Zelevinsky have developed an algorithm to realize the whole crystal for a given reduced expression for the longest element $w_{0}$ in the Weyl group. This algorithm has to do with taking a PBW datum with respect to the chosen reduced word, which becomes a vertex of the crystal, and applying the crystal operator, which becomes an edge of the crystal. Tingley, Salisbury, and Schultze found a simpler realization where the crystal operators are represented by a bracketing rule. This bracketing rule was restricted to one of the four classical Lie types, type $A_{n}$ as well as for a special case of word called a braidless word in the other three classical Lie types, $B_{n}, C_{n},$ and $D_{n}$.

We generalized the above bracket rule beyond braidless words for specific reduced expressions. This extends the use of this efficient algorithm above as it eliminates the issue of the restriction to braidless words only allowing for a bracket rule for the ordinary crystal operators or for the *-operators, which are reversed due to Kashiwara's involution, but not for both types of operators simultaneously. With our result, there is a simple combinatorial model of the algorithm in the case of both the ordinary crystal operators and the *-operators on a single realization of $B(\infty)$ in types $B_{n}, C_{n},$ and $D_{n}$.


This material is based upon work supported by the National Science Foundation under Grant No. DMS-1929284 while the authors were in residence at the Institute for Computational and Experimental Research in Mathematics in Providence, RI, during the Summer@ICERM 2026: Pure and Applied Mathematical Models program.
}
 \newpage